%%%%%%%%%%%%%%%%%%%%%%%%%%%%%%%%%%%%%%%%%%%%%%%%%%%%%%%%%%%%%%%%%%
%% Toto je inkludovaný LaTeXový soubor `c_conc.tex'.            %%
%% ------------------------------------------------------------ %%
%% Verze 2022-05-31                                             %%
%% Vyvíjí & udržuje <stanislav.hledik@physics.slu.cz>           %%
%%%%%%%%%%%%%%%%%%%%%%%%%%%%%%%%%%%%%%%%%%%%%%%%%%%%%%%%%%%%%%%%%%

\chapter*{Pár doporučení na závěr \ldots}\label{conclus}%##################

\ldots zejména pro studenty bakalářského studia.

Pro obhajobu závěrečné práce byste si měli vytvořit a dataprojektorem promítat
presentaci o~trvání cca 15 až 25~minut~-- to závisí na typu práce~-- a
rozhodně byste se o~délce trvání presentace měli pobavit s~vedoucím vaší
práce. Rozhodně při obhajobě nepromítejte stránky ze závěrečné práce vysázené
pomocí \pkg{\ipotname} (mějte je však pro potřeby případné diskuse
připraveny). Jednak díky portrétové orientaci špatně využívá promítací plochu
a má pro tyto účely příliš malý, obtížně čitelný stupeň písma, jednak by
presentace měla být pouze heslovitá, bez ucelených vět a odstavců. To, co je
vhodné k~tisku na papír, není vhodné pro presentaci a obráceně. Pro vytvoření
presentace použijte buď PowerPoint nebo některé z~\LaTeX
ových\INDEX{publikační systém!\LaTeX{}} řešení
\pkg{pdfscreen}~\cite{Rad:LIVE:CTAN:pdfscreen},
\pkg{Beamer}~\cite{samc:LIVE:CTAN:beamer}.

Při přípravě závěrečné práce pomocí \pkg{\ipotname} doporučuji vypnout volbu
\opt{htext} popsanou v~Dodatku~\ref{options} na stránce~\pageref{opthtext},
kompilace se tím urychlí. Kromě toho byste měli mít tuto volbu vypnutou
u~kompilace pro tištěnou formu práce, neboť zapnutá volba \opt{htext} barevně
odlišuje různé typy odkazů, a tato barva se bude tisknout\footnote{Na rozdíl
  od implicitních rámečků, které jsou vidět při prohlížení na monitoru, ale ne
  v~tisku.}  (v~případě černobílého tisku jako odpovídající stupeň
šedi). Naopak u~kompilace pro elektronickou verzi práce je vhodné mít volbu
\opt{htext} zapnutou; hypertextové odkazy urychlují prohlížení dokumentu.

V~názvech \LaTeX ových\INDEX{publikační systém!\LaTeX{}} zdrojových souborů a
souborů s~obrázky nepoužívejte akcentovaná písmena a mezery.  \TeX{} systém by
s~tím mohl mít potíže. Místo \verb!Třetění úhlu.tex! použijte
\verb!Treteni_uhlu.tex! a podobně.

Rozsáhlejší práce rozdělte na \emph{kořenový (mateřský) soubor} obsahující
standardní strukturu \LaTeX ového\INDEX{publikační systém!\LaTeX{}} dokumentu
(u~tohoto manuálu se jedná o~soubor \pgm{\thesnamecze .tex}), a~vlastní text
přesuňte (nejlépe po kapitolách) do separátních souborů, jenž vložíte do
kořenového souboru pomocí příkazů \verb!\include!. Podrobné instrukce
naleznete např.  v~\cite{Kop-Dal:2004:LaTeXPodrPruv:}. V~případě zdrojových
souborů tohoto manuálu jsou tyto inkludované soubory umístěny v~adresáři
\pgm{include}.\footnote{Dokumentní třída \pgm{\ipotname .cls} je v~něm najde.}

% Obsahuje-li závěr dokumentu jeden zcela prázdný list, je to v~pořádku.  Na
% konci práce by se totiž měla objevit vakát\INDEX{vakát} (prázdná stránka),
% takže končí-li práce sudou (levou) stránkou, musíme vložit \emph{dva}
% vakáty,\INDEX{vakát} tj.  prázdný list, ne jen jeden vakát.\INDEX{vakát}
% Dokumentní třída\INDEX{dokumentní třída} \file{\ipotname .cls}\INDEX{dokumentní
%   třída!\ipotname} to bez volby \opt{draft} při výstupu do \designation{dvi} dělá
% automaticky a není třeba se o~to starat (viz též Dodatek~\ref{necess},
% stránka~\pageref{necess}).

Máte-li v~práci enormní množství matematiky, stojí ze to podívat se do Wickovy
knihy~\cite{Wic:1966:PravMatSaz:} nebo do
článku~\cite{Hor:2001:ZprCSTUG:MathTypeset} a nastudovat základní pravidla
matematické sazby. Jako příručka matematické sazby může posloužit
i~text~\cite{Nec:2003::YetiTypoBest}. Obecné pravidlo zní: nevíte-li si se
sazbou nějakého prvku rady, napište si
makro~\cite{Kop-Dal:2004:LaTeXPodrPruv:} s~prozatímní definicí, kterou můžete
po konzultaci nebo nastudování literatury dodatečně změnit. Pro obecné poučení
o~estetické sazbě stoji za to nahlédnout do~\cite{Pop-Fle-Pop:1984:SazbaI:}
nebo do Beranova \emph{Typografického manuálu}~\cite{Ber:1999:TypoMan:}.

\endinput

%%
%% End of file `c_conc.tex'.
